% =====================================================================
%  CHAPITRE 1 - PRÉSENTATION DE L'ORGANISME D'ACCUEIL
% =====================================================================
%  NOTES (à traiter avant version finale) :
%  [1] Chiffres publics : à revérifier (voir references.bib).
%  [2] Dénomination du pôle "contracting" : à confirmer sur l'URD 2024.
%  [3] Historique GTIE : formulation en commentaire, à activer après vérif.
%  [4] Liste des directions de VESI : laissée générique (non confirmée).
%  [5] Données internes : créditées "documentation interne VESI".
% =====================================================================

\chapter{Présentation de l'organisme d'accueil}
\label{chap:organisme}

Ce chapitre présente le cadre dans lequel s'est déroulé mon stage, du plus
général au plus spécifique~: le groupe VINCI, la branche VINCI~Energies,
son entité informatique VESI, la direction Custom Applications, puis le
service Collaborative Solutions et l'équipe qui m'a accueilli. Il se termine
sur l'engagement sociétal et environnemental du groupe et sa traduction
concrète dans les pratiques de développement.

\section{Le groupe VINCI}
\label{sec:groupe-vinci}

VINCI est un acteur mondial des concessions, de l'énergie et de la
construction, présent dans plus de 120~pays. Le groupe intervient sur
l'ensemble du cycle de vie des infrastructures et des bâtiments, de la
conception à l'exploitation, et inscrit son action dans une stratégie de
transition énergétique et environnementale. En 2024, il réunissait environ
285\,000~collaborateurs pour un chiffre d'affaires de 71,6~milliards d'euros
et un résultat net part du Groupe de 4,86~milliards d'euros~\cite{vinci_urd2024}.

Le groupe organise ses activités autour de deux grands pôles opérationnels.
\textbf{VINCI~Concessions} conçoit, finance et exploite des infrastructures
de transport, notamment à travers VINCI~Autoroutes et VINCI~Airports. Le pôle
\textbf{énergie et construction} rassemble \textbf{VINCI~Energies},
\textbf{VINCI~Construction} et \textbf{Cobra~IS}. À ces activités s'ajoute
\textbf{VINCI~Immobilier}, dédié à la promotion immobilière. C'est au sein de
VINCI~Energies que s'est déroulé mon stage.

\begin{figure}[H]
  \centering
  \begin{tikzpicture}[
    font=\footnotesize,
    racine/.style={rectangle, rounded corners, draw=vinciblue, very thick,
      fill=vinciblue, text=white, minimum width=2.4cm, minimum height=0.8cm,
      align=center},
    pole/.style={rectangle, rounded corners, draw=vinciblue, thick,
      fill=vinciblue!10, minimum width=2.6cm, minimum height=0.75cm, align=center},
    feuille/.style={rectangle, rounded corners, draw=gray, thick,
      fill=grisclair, minimum width=2.4cm, minimum height=0.65cm, align=center},
    lien/.style={-, draw=gray, thick}
  ]
    \node[racine] (vinci) {VINCI};
    \node[pole, below left=1.2cm and 3.1cm of vinci] (conc) {VINCI\\Concessions};
    \node[pole, below=1.2cm of vinci] (contr) {Énergie \&\\Construction};
    \node[pole, below right=1.2cm and 3.1cm of vinci] (immo) {VINCI\\Immobilier};
    \draw[lien] (vinci) -- (conc);
    \draw[lien] (vinci) -- (contr);
    \draw[lien] (vinci) -- (immo);
    \node[feuille, below=1.9cm of contr] (energies) {VINCI\\\textbf{Energies}};
    \node[feuille, left=0.35cm of energies] (construction) {VINCI\\Construction};
    \node[feuille, right=0.35cm of energies] (cobra) {Cobra~IS};
    \draw[lien] (contr) -- (construction);
    \draw[lien] (contr) -- (energies);
    \draw[lien] (contr) -- (cobra);
    \node[feuille, below=1.9cm of conc, minimum width=2.1cm] (airp) {VINCI\\Airports};
    \node[feuille, left=0.35cm of airp, minimum width=2.1cm] (auto) {VINCI\\Autoroutes};
    \draw[lien] (conc) -- (auto);
    \draw[lien] (conc) -- (airp);
    \begin{scope}[on background layer]
      \node[draw=vincired, thick, rounded corners, fit=(energies), inner sep=3pt] {};
    \end{scope}
  \end{tikzpicture}
  \caption{Structure simplifiée du groupe VINCI. L'entité d'accueil,
    VINCI~Energies, est encadrée en rouge.}
  \label{fig:structure-vinci}
\end{figure}

\section{La branche VINCI Energies}
\label{sec:vinci-energies}

VINCI~Energies accompagne ses clients dans la transition énergétique et la
transformation numérique, en proposant des solutions multitechniques sur
mesure, de la conception-réalisation à l'exploitation-maintenance. En 2024,
la branche a réalisé un chiffre d'affaires de 20,4~milliards d'euros, avec
plus de 100\,000~collaborateurs répartis dans 61~pays~\cite{ve_chiffres_cles}.

Sa principale caractéristique tient à son \textbf{modèle décentralisé}~:
VINCI~Energies fonctionne comme un réseau de près de 2\,000~entreprises à
taille humaine, ancrées dans leurs territoires, ce qui conjugue réactivité
locale et puissance d'un grand groupe. Ses métiers se répartissent entre
infrastructures d'énergie, industrie, technologies de l'information et
bâtiment, portés par des marques transverses telles qu'Actemium, Axians,
Omexom, Citeos et VINCI~Facilities. Ce modèle a une conséquence directe pour
les systèmes d'information~: outiller un grand nombre d'entités hétérogènes
avec des applications communes, fiables et sécurisées. C'est précisément la
mission de VESI, et le contexte de mon stage.

% [3] HISTORIQUE GTIE - à activer seulement après vérification de source :
% VINCI~Energies s'inscrit dans la continuité d'un savoir-faire ancien en
% génie électrique, hérité notamment de l'entité GTIE.

\begin{table}[H]
  \centering
  \caption{Carte d'identité de VINCI~Energies (données 2024).}
  \label{tab:carte-ve}
  \begin{tabularx}{\textwidth}{>{\bfseries}p{4.2cm} X}
    \toprule
    \normalfont\bfseries Catégorie & \bfseries Informations \\
    \midrule
    Raison sociale        & VINCI~Energies SA \\
    Métiers               & Services à l'énergie, industrie, ICT, bâtiment \\
    Présence              & 61 pays, près de 2\,000 entreprises \\
    Effectifs (2024)      & plus de 100\,000 collaborateurs \\
    Chiffre d'affaires (2024) & 20,4 milliards d'euros \\
    Marques               & Actemium, Axians, Omexom, Citeos, VINCI~Facilities \\
    Concurrents           & Equans, Spie, Eiffage Énergie Systèmes, SNEF \\
    \bottomrule
  \end{tabularx}
\end{table}

\section{VINCI Energies Systèmes d'Information (VESI)}
\label{sec:vesi}

VINCI~Energies Systèmes d'Information (VESI) est l'entité chargée du système
d'information de VINCI~Energies. Sa mission est de concevoir, déployer et
maintenir les applications et outils numériques utilisés par les entreprises
du réseau, qu'il s'agisse d'outils de gestion, de collaboration ou de
valorisation de la donnée. Ce rôle est stratégique dans un groupe
décentralisé~: fournir un socle applicatif commun à des centaines d'entités
autonomes suppose des solutions robustes, sécurisées et adaptables. VESI
dispose d'une implantation notable au Mans, où s'est déroulé mon stage. Elle
regroupe plusieurs directions spécialisées, parmi lesquelles la direction
\textbf{Custom Applications}, qui m'a accueilli. La
figure~\ref{fig:org-vesi}, en fin de chapitre, restitue l'ensemble de cette
chaîne organisationnelle.

\section{La direction Custom Applications}
\label{sec:custom-apps}

La direction Custom Applications conçoit, développe, teste, déploie et
maintient les applications métier --- web et mobiles --- de VINCI~Energies
(hors périmètres Codex, FSM et CRM). Elle maintient un portefeuille de plus
de 50~applications couvrant des besoins variés~: outils Groupe, applications
opérationnelles, QHSE, solutions collaboratives et services internes à
VESI\footnote{Source~: documentation interne VINCI~Energies Systèmes
d'Information.}.

Ses développements reposent sur un socle full-stack --- C\#/.NET, Angular,
React, Java, Flutter et technologies mobiles --- complété par le déploiement
et le support de solutions éditeurs (paie, SIRH). Au-delà du développement,
la direction assure la gouvernance des solutions, l'animation des communautés
d'utilisateurs clés et intègre les pratiques \textbf{DevSecOps} et de
\textbf{numérique responsable} dans ses projets. Elle représentait environ
98~collaborateurs (ETP) en 2025, sur un périmètre France et
international\footnote{Source~: documentation interne VINCI~Energies Systèmes
d'Information.}.

Pour structurer cette activité, la direction est organisée en quatre services,
chacun responsable d'un domaine applicatif~: \textbf{Web \& Mobile
Applications}, \textbf{Collaborative Solutions}, \textbf{Payroll Solutions}
et \textbf{HRIS}. J'ai été intégré au service Collaborative Solutions.

\section{Le service Collaborative Solutions}
\label{sec:collab-solutions}

Le service Collaborative Solutions maintient et développe les solutions
collaboratives de VINCI~Energies, principalement sur une base
\textbf{Microsoft~365}, dans le respect des règles de gouvernance et de
sécurité. Son rôle est transverse~: renforcer la collaboration au sein d'un
groupe fortement décentralisé, en fournissant des outils communs à des
centaines d'entités.

Le service s'articule autour de deux équipes\footnote{Source~: documentation
interne VINCI~Energies Systèmes d'Information.}. L'équipe \textbf{Group
Collaborative Solutions} prend en charge les solutions à périmètre Groupe ---
notamment l'intranet \textbf{MyVIEW}, le moteur de recherche et annuaire
\textbf{Connect}, et les espaces de travail collaboratifs. L'équipe
\textbf{BU Collaborative Solutions} se concentre sur les solutions à périmètre
Business Unit, comme \textbf{DPF} (Digital Project Files), ainsi que sur les
migrations de données M365 (OneDrive, SharePoint, Teams). Sur le plan
technique, les compétences clés du service couvrent le développement
C\#/.NET, Angular et React, appuyé sur SharePoint Online et l'API
Microsoft~Graph.

C'est au sein de l'équipe \textbf{Group Collaborative Solutions} que s'est
déroulé mon stage, plus précisément dans le squad M365, l'un des deux squads
qui la composent avec le squad MyVIEW. Mes missions, détaillées dans les
chapitres suivants, se sont inscrites dans ce périmètre --- au premier rang
desquelles la refonte de l'application \csbag. La figure~\ref{fig:org-vesi}
récapitule l'ensemble de la chaîne organisationnelle, de VESI jusqu'au squad
d'accueil.

\begin{figure}[H]
  \centering
  \begin{tikzpicture}[
    font=\scriptsize,
    racine/.style={rectangle, rounded corners=2pt, draw=vinciblue, very thick,
      fill=vinciblue, text=white, minimum width=2.6cm, minimum height=0.7cm,
      align=center},
    boite/.style={rectangle, rounded corners=2pt, draw=gray, thick,
      fill=grisclair, minimum height=0.8cm, align=center, inner sep=3pt},
    hi/.style={rectangle, rounded corners=2pt, draw=vincired, very thick,
      fill=vinciblue!10, minimum height=0.8cm, align=center, inner sep=3pt},
    lien/.style={-, draw=gray, thick}
  ]
    \node[racine] (vesi) at (-2.5,0) {VESI};

    \node[hi, text width=2.2cm] (ca) at (-2.5,-1.5) {Direction\\Custom Applications};
    \node[boite, text width=2.2cm] (autres) at (1.6,-1.5) {Autres\\directions};

    \node[boite, text width=2cm] (s1) at (-6.7,-3.3) {Web \& Mobile\\Applications};
    \node[hi, text width=2cm]    (s2) at (-3.9,-3.3) {Collaborative\\Solutions};
    \node[boite, text width=2cm] (s3) at (-1.1,-3.3) {Payroll\\Solutions};
    \node[boite, text width=2cm] (s4) at (1.7,-3.3)  {HRIS};

    \node[hi, text width=2.9cm]    (grp) at (-5.7,-5.1) {Group\\Collaborative Solutions};
    \node[boite, text width=2.9cm] (bu)  at (-2.1,-5.1) {BU\\Collaborative Solutions};

    \node[hi, text width=2.1cm]    (m365) at (-7.2,-6.8) {Squad\\M365};
    \node[boite, text width=2.1cm] (myv)  at (-4.2,-6.8) {Squad\\MyVIEW};

    \draw[lien] (vesi) -- (ca);
    \draw[lien] (vesi) -| (autres);
    \foreach \s in {s1,s2,s3,s4} \draw[lien] (ca.south) -- (\s.north);
    \draw[lien] (s2.south) -- (grp.north);
    \draw[lien] (s2.south) -- (bu.north);
    \draw[lien] (grp.south) -- (m365.north);
    \draw[lien] (grp.south) -- (myv.north);
  \end{tikzpicture}
  \caption{Chaîne organisationnelle de VESI au squad d'accueil. Les entités
    encadrées en rouge situent le parcours menant au squad M365.}
  \label{fig:org-vesi}
\end{figure}

\section{Responsabilité sociétale et environnementale (RSE)}
\label{sec:rse}

La responsabilité sociétale et environnementale occupe une place centrale
dans la stratégie de VINCI~Energies, qui s'engage à réduire son empreinte
carbone et à accompagner ses clients dans leur transition énergétique. Cet
engagement se prolonge au niveau de VESI et de la direction Custom
Applications à travers deux axes intégrés directement aux développements~: la
\textbf{sécurité dès la conception} (approche DevSecOps, analyse systématique
du code et des dépendances) et le \textbf{numérique responsable}
(éco-conception des applications, accessibilité
numérique)\footnote{Source~: documentation interne VINCI~Energies Systèmes
d'Information.}. Les enjeux RSE ne sont donc pas seulement portés au niveau du
groupe~: ils irriguent aussi les pratiques quotidiennes des équipes de
développement.